\scanpage{565}
\mspage{476}
de a) ci-dessus, quand $C'$ est \uline{sans projecteurs}.

\uline{Question} 4.16.6.4. Les conditions 1°) et 2°) du
\unsure{précédentes|précédent} impliquent-elles celles de 4.16.6.2 ?
En d'autres termes, si la catégorie base
$C$ est égale à $\unsure{Q|\mathcal Q}$ ou à $Q'$, est-ce que ces
conditions impliquent-elles que $\unsure{\mathcal E|E}$ est
accessible, et $f:\mathcal E\longrightarrow\deleted{\unsure{\mathcal E}}C$ accessible ?

\unsure{Ça|Sa} semble \unsure{idiot} de ne pas connaître
la réponse !

\S 4.16.6.5. J'ai envie d'appeler
catégorie \uline{accessible sur} $\deleted{\unsure{\mathcal E}}C$, une catégorie
$\mathcal E$ sur $C$ satisfaisant les conditions
équivalentes de \S 4.16.6.1 --- du moins
dans le cas où on suppose $C$ petite.
Pour que cette terminologie soit vraiment
\unsure{satisfaisante|satisfaisant}, il faudrait que l'on
sache prouver que ces conditions impliquent
ceci : pour toute catégorie accessible
$M$ et \unsure{tout} foncteur accessible $M\longrightarrow C$,
la catégorie $\deleted{\unsure{M|C}}\mathcal E\times_C M$ est accessible, et
\unsure{les|le} \unsure{foncteurs|foncteur} $\mathcal E\times_C M\mathrel{\unsure{-|\longrightarrow}}M$ est accessible.
% archive sheet number: 564
